\chapter{Joint Pain}

\section*{Definition}
Pain or discomfort localized to one or more joints varying from dull ache to sharp stabbing sensation.

\section*{What Happens in Your Body}
The number and pattern of affected joints—single or multiple, symmetric or asymmetric, inflammatory or degenerative—help guide the differential diagnosis.

\section*{Causes}
\begin{itemize}
  \item Osteoarthritis (degenerative)
  \item Rheumatoid arthritis (inflammatory)
  \item Trauma/injury
  \item Septic arthritis (infectious)
  \item Gout (crystal arthropathy)
\end{itemize}

\section*{When to See a Doctor}
See your doctor if:
\begin{itemize}
  \item Joint pain is persistent, severe, or accompanied by swelling or redness
\end{itemize}

\section*{Related Symptoms}
\begin{itemize}
  \item Swelling
  \item Warmth/redness
  \item Morning stiffness >30min
  \item Reduced movement
  \item Systemic symptoms like fever/rash
\end{itemize}

\section*{Self-Care / Home Management}
\begin{itemize}
  \item Use relative rest, ice for a recent injury, and gentle movement within comfort; avoid prolonged immobilization.
  \item Acetaminophen or an NSAID may help when safe and used according to the label.
  \item A hot, swollen joint with fever or sudden severe pain needs urgent assessment rather than home treatment.
\end{itemize}

\section*{How Your Doctor Will Evaluate You}
\begin{itemize}
  \item History and examination assess the number of joints, symmetry, inflammatory features, trauma, fever, rash, and extra-articular symptoms.
  \item Blood tests and imaging are selected according to the clinical pattern; aspiration of an acutely swollen joint may be needed to exclude infection or crystals.
\end{itemize}

\section*{Treatment Options}
\begin{itemize}
  \item Treatment may include exercise and physical therapy, analgesia, joint aspiration or injection, disease-modifying therapy, or antibiotics for septic arthritis.
  \item Persistent inflammatory symptoms warrant timely rheumatology assessment; structural injury may require orthopaedic care.
\end{itemize}

\section*{References}
\begin{thebibliography}{99}
\bibitem{joint_pain:ref1} Zhang W, Doherty M, Peat G, et al. ``EULAR evidence-based recommendations for the diagnosis of knee osteoarthritis.'' \textit{Annals of the Rheumatic Diseases}, 2010;69(3):483--489.
\bibitem{joint_pain:ref2} Aletaha D, Neogi T, Silman AJ, et al. ``2010 Rheumatoid arthritis classification criteria.'' \textit{Arthritis \& Rheumatism}, 2010;62(9):2569--2581.
\bibitem{joint_pain:ref3} Neogi T, Jansen TLTA, Dalbeth N, et al. ``2015 Gout classification criteria.'' \textit{Arthritis \& Rheumatology}, 2015;67(10):2557--2568.
\bibitem{joint_pain:ref4} Mathews CJ, Weston VC, Jones A, et al. ``Bacterial septic arthritis in adults.'' \textit{Lancet}, 2010;375(9717):846--855.
\bibitem{joint_pain:ref5} Felson DT. ``Osteoarthritis of the knee.'' \textit{New England Journal of Medicine}, 2006;354(8):841--848.
\bibitem{joint_pain:ref6} Dieppe PA, Lohmander LS. ``Pathogenesis and management of pain in osteoarthritis.'' \textit{Lancet}, 2005;365(9463):965--973.

\end{thebibliography}
